Hemos demostrado a través del presente documento empírico que la maximización del valor de inteligencia predictiva no rige de forma exclusiva en la precisión estadística, sino fundamentalmente en su nivel de explicabilidad y acoplamiento directo con la estructura de costos empresariales reales. 

El desarrollo guiado longitudinalmente estableció una arquitectura fluida donde los datos transaccionan desde sus orígenes en bruto, se destilan computacionalmente mediante potentes ensambles de árboles con decodificación humana XAI (SHAP), y culminan de forma intuitiva como tableros interactivos para los niveles directivos. Este trabajo ha trasladado exitosamente los tradicionales prototipos puros orientados a la academia e integró las perspectivas de Inteligencia de Negocios para prescribir rentabilidad.

Estudios subsecuentes deberían apuntar al levantamiento de diseños causales experimentales (A/B testing iterativos controlados) que comprueben que las palancas expuestas como factores causadores de la probabilidad final por los algoritmos reaccionen, una vez operadas activamente por los operadores de servicio, verdaderamente mitigando la huida prospectiva.
